%%%%%%%%%%%%%% LAYOUT & FONTS %%%%%%%%%%%%%%%%%%%%%

\documentclass[12pt]{article}
\usepackage {setspace}
\onehalfspacing
\pagestyle{plain}
\usepackage{fontspec}% Enable KEYMAN functionality
\setmainfont{Libertinus Serif}
%\usepackage{indentfirst}
\setlength{\parindent}{1cm}
\usepackage{authblk}
\usepackage{avant}
\usepackage[width=155mm,top=25mm,bottom=25mm]{geometry}
\usepackage{booktabs}
\usepackage{capt-of}
\usepackage[labelfont=bf]{caption}
\usepackage{graphicx}
\usepackage{media9}
\usepackage{lineno}
\linenumbers
\usepackage[bottom]{footmisc}
%\tracingmacros=1
%\tracingcommands=2

%%%%%%%%%%%%%%%TABLLES%%%%%%%%%%%%%%%%%%%%%%
\usepackage{ltablex} % Solutions here. https://tex.stackexchange.com/questions/375024/table-number-not-right-after-using-tabularx
\usepackage{longtable}
%\usepackage{makecell}
%\newcolumntype{Z}{ >{\raggedright\arraybackslash}X }\usepackage{capt-of}
\usepackage{tocloft, calc} % calc is added to have "Chapter" in ToC
%\usepackage[titles]{tocloft} % Include [titles] in case you want to change the size of ToC, LoT, etc
%\renewcommand{\cfttoctitlefont}{\hfill\large\bfseries}
%\renewcommand{\cftloftitlefont}{\hfill\large\bfseries}
%\renewcommand{\cftlottitlefont}{\hfill\large\bfseries}
%\usepackage{calc}
%\usepackage[labelfont=bf]{caption}%%%Comment out for bolded Caption
\usepackage{float}
\usepackage{caption}
\usepackage{subcaption}
\usepackage{multirow}
\usepackage{multicol}
\usepackage{tablefootnote}
%\setlength{\tabcolsep}{1pt} % default is 6pt for column width
 
%%%%%%%REFERENCES%%%%%%%%%%%%%%%%% 
\usepackage[sort]{natbib}
\bibpunct[: ]{(}{)}{,}{a}{}{,}
\setlength{\bibsep}{0.0pt}


%%%%%%%LINKS AND HYPERLINKS%%%%%%%%%%%%%%%%
\usepackage[hidelinks,hypertexnames=false]{hyperref}
\hypersetup{colorlinks=true,
    linkcolor=blue,
    filecolor=blue,      
    urlcolor=blue,
    citecolor=blue}
\urlstyle{same}
\usepackage{bookmark}
\bookmarksetup{numbered} %solution for numbered bookmark sections

%%%%%%%%%%%%Glosses & related packages%%%%%%%%%%%%%%%
%%%\usepackage[inline]{leipzig}
\usepackage[mcolblock,leipzighyper]{leipzig}
\let\eachwordone=\itshape % italic data line (default is \rm)
\glssetwidest{xxxxxxx}
\usepackage{gb4e}
\noautomath	
\renewcommand{\leipzigname}{Abbreviations and special symbols}
%glosses.tex
%%%%%%%%%%% MOCOVI GLOSSES%%%%%%%%%%%%%%%%%%%%%%%%%%%%%%%%%%%%
%%%%%%%%Special symbols%%%%%%%%
\newleipzig{sploan}{$_{sp}$}{Spanish loanword}
\newleipzig{tbloan}{$_{tb}$}{Toba loanword}
\newleipzig{ungr}{$^{*}$}{ungrammatical}
\newleipzig{clitic}{=}{clitic ligature}
\newleipzig{dirch}{$\rightarrow$}{direction of change}

%%%%%%%Specyfic glosses%%%%%%%%%%%%%%%%%%%%%%%%%%%%%%%%%%%%%%%%%
\newleipzig{coord}{coord}{coordinant}
\newleipzig{link}{link}{linker}
\newleipzig{disc}{disc}{discourse maker}
\newleipzig{tprl}{tprl}{temporal}
\newleipzig{evd}{evd}{evidential}
\newleipzig{trvz}{trvz}{transitivizer} 
\newleipzig{III}{III}{Set III bound person form}
\newleipzig{II}{II}{Set II bound person form}
\newleipzig{I}{I}{Set I bound person form}
\newleipzig{dir}{dir}{directional}
\newleipzig{arg}{arg}{argument}
\newleipzig{defc}{defc}{defocusing}
\newleipzig{dim}{dim}{diminutive}
\newleipzig{class}{class}{classifier}
\newleipzig{act}{act}{action}
\newleipzig{vm}{vm}{valence modifier}
\newleipzig{exist}{exist}{existential}
\newleipzig{noml}{noml}{nominal}
\newleipzig{ints}{ints}{intensifier} 
\newleipzig{rec}{rec}{recipient}
\newleipzig{causee}{causee}{causee argument}
\newleipzig{coll}{coll}{collective}
\newleipzig{pron}{pron}{pronominal}
\newleipzig{des}{des}{desiderative}
\newleipzig{val}{val}{valence}
\newleipzig{exclam}{exclam}{exclamation}
\newleipzig{inter}{inter}{interrogative}
\newleipzig{vrbz}{vrbz}{verbalizer}
\newleipzig{iter}{iter}{iterative}
\newleipzig{attr}{attr}{attributive}
\newleipzig{ac}{ac}{active}% Grondona gloss
\newleipzig{indet}{indet}{indeterminate}%
%%%%%%%%%%%%LOCATIVES%%%%%%%%%%%%%%%%%%%%%%%%
\newleipzig{locOne}{loc$_{1}$}{locative type 1}%lek
\newleipzig{locTwo}{loc$_{2}$}{locative type 2}%ge
\newleipzig{locThree}{loc$_{3}$}{locative type 3}%git
\newleipzig{locFour}{loc$_{4}$}{locative type 4}%gi
%%%%%%%%%%DIRECTIONALS
\newleipzig{dirDown}{dir$_{down}$}{directional: downwards}%
\newleipzig{dirOut}{dir$_{out}$}{directional: outwards}
\newleipzig{dirIn}{dir$_{in}$}{directional: inwards}
\newleipzig{dirUp}{dir$_{up}$}{directional: upwards}
\newleipzig{out}{out}{outwards}%wek
\newleipzig{down}{down}{downwards}%ɲi
\newleipzig{up}{up}{upwards}%ʃigem 
\newleipzig{in}{in}{inwards}%ʃigem 

% %%%%%%%%%DETERMINERS%%%%%%%%%%%%%%%%%%%%%%%%
\newleipzig{detOne}{det$_{1}$}{determiner type 1: non-visible} %ka
\newleipzig{detTwo}{det$_{2}$}{determiner type 2: going, far}%so
\newleipzig{detThree}{det$_{3}$}{determiner type 3: coming, close}%na
\newleipzig{detFour}{det$_{4}$}{determiner type 4: static, vertical}%ɾa
\newleipzig{detFive}{det$_{5}$}{determiner type 5: static, horizontal} %dʒi
\newleipzig{detSix}{det$_{6}$}{determiner type 6: static, sitting }%ɲi
%\newleipzig{deicg}{deic$_{(gng)}$}{deictic root: going} %Grondona gloss
%%%%%%%%%%%TODO: Replace these glosses%%%%%%%%%%%%%%%%
%\newleipzig{dsit}{dsit}{determiner sitting}%ɲi
%\newleipzig{dhor}{dhor}{determiner horizontal}%dʒi
%\newleipzig{dstd}{dstd}{determiner standing}%ɾa
%\newleipzig{dgng}{dgng}{determiner going,far}%so1
%\newleipzig{dnvs}{dnvs}{determiner non-visible, away}%ka
%\newleipzig{dccl}{dccl}{determiner coming, close}%na

\makeglossaries
%%%%%%%%%%%%%AUTHOR-TITLE%%%%%%%%%%%%%%%5
\title{\textbf{Agent foregrounding and the conditioning of P‑argument expression: qualitative and quantitative observations}}
\author{}
%\affil{The University of Texas at Austin\vspace{-6mm}}
%\affil{\underline{cristianrj@utexas.edu}}
\date{Version \today}


%%%%%%%%%%%%%%DOCUMENT%%%%%%%%%%%%%%%%%%%%%%
\begin{document}
\maketitle

\begin{abstract}
This paper examines agent foregrounding constructions in Northern Chaco Mocoví and their effects on the expression of P arguments. Building on a description of clause types and argument expression, the analysis focuses on the derivational markers \textit{‑ɢan} and \textit{‑taɢan}, which combine detransitivizing and causativizing functions. When used as intransitivizers, these markers increase the prominence of the agent and demote the P argument, giving rise to a set of constructions functionally comparable to antipassives. Based on both elicited and natural discourse data, the study identifies the morphosyntactic correlates of P‑argument demotion, the conditions under which P arguments may or may not be expressed, and the discourse strategies used to reintroduce them. The paper also shows that the detransitivized forms play a central role in lexicon building, especially in the formation of derived nouns. A quantitative case study complements the qualitative analysis by illustrating the distribution of agent foregrounding  and morphologically related constructions in natural speech. The findings contribute to understanding valence alternation, argument expression, and lexical derivation in Mocoví and to broader typological discussions on antipassive‑related phenomena.
\end{abstract}

\noindent\textit{Keywords}: agent foregrounding, P-demotion, Northern Chaco Mocoví, lexical distribution, frequency 

\section{Introduction}

This paper examines the morphosyntactic and discourse properties of agent foregrounding constructions in Northern Chaco Mocoví. These constructions arise through the use of the derivational markers \textit{‑ɢan} and \textit{‑taɢan}, which participate in a broader system of valence‑changing morphology. When used as detransitivizers, these markers reduce the verbal valence of two‑place predicates and highlight the agent’s activity, contributing in the demotion or omission of the P argument. Because this function overlaps with properties commonly associated with antipassive constructions, the analysis also draws on recent typological work on antipassives to contextualize the Mocoví patterns.

The study first describes the relevant clause types and the mechanisms by which Mocoví expresses core arguments. Particular attention is given to the role of nominal determiners, the interaction between verb morphology and argument structure, and the conditions that license or block the expression of the P argument. The paper then explores how the detransitivizing use of \textit{‑ɢan} and \textit{‑taɢan} affects the presence, form, and interpretation of the P argument, especially in relation to definiteness, referentiality, and argument indexing.

Based on elicited data and natural discourse, the analysis identifies several strategies for recovering or reintroducing the P argument after demotion, including multiverbal sequences and oblique‑marked participants. The distribution of these strategies is examined in relation to verb valence, predicate semantics, and narrative context. Additionally, the paper shows that detransitivized predicates play an important role in lexicon building, serving as bases for derived nouns formed through the nominalizer \textit{‑ɢak}.

A small case‑study based on a narrative text provides a first quantitative perspective on the frequency and discourse behavior of these constructions. Although agent foregrounding predicates occur relatively infrequently in discourse, their distribution is systematic and reflects functional motivations related to participant tracking and event structure. Overall, the study contributes to understanding how Mocoví organizes valence alternation, argument prominence, and lexical derivation within its verbal system.

%Languages of the Gran Chaco region, particularly those of the Guaycuruan family, exhibit a noteworthy valence alternation pattern in which a single derivational marker encodes both antipassive‑like and causative functions. Although well known to specialists working on Mocoví, Toba, and related languages \citep[see, among many other][]{carpio2012a,Censabella2008,censabella2002,Gonzalez2015,gualdieri1998,carrio2009}, this co‑expression pattern remains under‑discussed in the broader typological literature. Previous studies have highlighted that the antipassive‑like function of these markers is associated with an increase in the agentive prominence of the subject argument, resulting in the backgrounding or demotion of the erstwhile P argument.

%This paper investigates agent foregrounding constructions in Northern Chaco Mocoví and the consequences of this process for P‑argument expression in discourse. I adopt agent foregrounding as a language‑internal descriptive label that integrates into the broader system of valence‑changing morphology in Mocoví. The study combines qualitative and quantitative methods: detailed morphosyntactic analysis of elicited and naturalistic data, and a small corpus‑based examination of the discourse frequency and distribution of derived predicates.

%The goals of this work are fourfold: (i) to describe the mechanisms by which P arguments may be demoted; (ii) to document the available structural strategies for expressing or omitting P arguments in agent foregrounding contexts; (iii) to assess the relative frequency and discourse behavior of these constructions; and (iv) to explore their implications for lexical derivation and lexicon building. By situating Mocoví within larger typological discussions on antipassives, valence changes, and derivational morphology, the paper contributes to understanding the structural diversity and functional motivations underlying voice‑related alternations across languages.

%A group of South American languages from the Gran Chaco region, belonging to the Guaycuruan family, possess a set of derivational verb markers co-expressing two opposite functions, closely related to the domain of antipassive and causative. While this co-expression pattern is well-know for scholars working on Guaycuruan and neighboring Gran Chaco languages , little is know about this phenomenon outside the South American context. Scholars working on those derivational markers show differences on the way that the definitional parameters related to these markers are weighed. There is, however, a recurrent observation on that the valency-decreasing, antipassive-like function primarily serves to indicate the agentivity salience and increase of the subject argument. Such a view comes up, for example, in the different glosses and segmentation proposed for cognate forms in Guaycuruan languages. A by-product of the increase in agentivity of the subject argument is the backgrounding of the former object argument. This paper explores agent foregrounding and its conditioning of object argument expression in Mocoví, a Guaycuruan language spoken in northeastern Argentina. I follow the term agent foregrounding as a language-specific label, which fits and coherently integrates to the more broad derivational system of valence-changing markers in Mocoví (and Guaycuruan). This work combines a qualitative and quantitave analysis of agent foregrounding constructions.

%Goals
%\begin{itemize}
%\item Ways of demoting the object: explain the derivational mechanisms and show the coherence on subject orientation. 
%\item Possible expression types.
%\item Frequency at the discourse level and correlations with properties observed for similar constructions in the domain of antipassives. 
%\item Implications of the derivational mechanism for lexicon building and other strategies at the word level of lexical derivation. Point to some comparative aspects with Mataguayan languages \citep[see][]{Nercesian2015b}.
%\end{itemize}


   

%\section{Family overview and glossing}
%Remember that call of complex affixes. Some ideas are important from \citet{dom.etal2015}. Some of the functions her resemble some of the features that linked to antipassive and the correlations of antipassive markers to other domain. Specifically, the notion of 

\section{The Mocoví people and their language}
Mocoví is part of the Guaycuruan family along with other four languages, three of which are still spoken today, i.e., Kadiwéo, Toba and Pilagá. Mocoví communities are located in northeastern Argentina, as shown in Map \ref{fig:moc-communities}\footnote{Personal elaboration based on the name of Mocoví communities as reported in different studies on Mocoví language and culture; see among others, \citet[][]{Altman2010,Altman2011}, \citet{carrio2009}, \cite{Dalla-CorteCaballero2012}, \citet{gualdieri1998}, \citet{grondona1998}, \citet{iparraguirre2011,Iparraguirre2016}, \citet{Lopez2017}.} They are spread out across two neighboring provinces, Chaco and Santa Fe, creating a chain of geographically connected clusters of Mocoví families.

\begin{figure}[H] 
\centering 
    \includegraphics[width=1\linewidth]{/Users/cristian/Documents/repos/maps/mocovi-argentina/plot/moc-argentina.png}
    \caption{Mocoví communities in the Argentine Gran Chaco.}
    \label{fig:moc-communities}
\end{figure}		



Demographically, as one moves further north within the chain of Mocoví communities, it becomes increasingly common  to find Mocoví people who have been in contact with Toba-speaking communities for over a century or more. This is the case, for example, in the community Cacica Dominga, located in the larger rural area of Colonia Aborigen, which marks one of the northernmost points of the Mocoví community chain in Argentina. Here, Mocoví and Toba families have been in a long-term and intense contact, fostering a complex social organization that also integrates non-indigenous people. The specific social and linguistic outcomes of this long-term contact situation are still in need of further research.

Data included in this study mostly come from what I call Northern Chaco Mocoví, which corresponds to the language as spoken in Cacica Dominga and closely connected neighbouring settlements in the southern region of Colonia Aborigen \citep[see][]{Juarez2019}. The large corpora included here deals with controlled and natural speech data collected originally via fieldwork. Northern Chaco Mocoví data include a basic metadata information indicating date of collection and location of the example in the audio-video files e.g., (mocCA160720\_1: 00:31:19). When examples come from a source other than my own collection, this information is also provided in parenthesis for the specific illustration. See also \autoref{sec:supplementary} to know more about the accompanying documentation of this paper, i.e., repository and access to raw data. 


\section{Morphosyntactic language profile}\label{sec:morphosynstactic-overview}

Mocoví, as well as it sister Guaycuruan languages, is a head-marking language, with a quite robust verb/noun distinction \citep[see, for instance][]{juarez2022a}. Verbs and nouns show a complex morphological structure, where nominal morphology distinguishes obligatorily possessed, optionally possessed and non-possessed nouns. Additionally, the language lacks case-marking and possess a single oblique marker for the nominal domain.  The verb complex, on the other hand, includes dedicated morphological elements expressing person and number of core arguments, spatial information such as location and direction, aspect, change of valence, negation and source of information/event temporality. Consider the initial examples in (\ref{ex:determinerplural}) and (\ref{ex:verbstructure1}). 
 
\begin{exe}
\ex 
\gll   so-wa jale-ɾ i-walqa-taʔ-pe ke-ɾa ʔoːtʃi \\ 
{\DetTwo}-{\Pl} man-{\Pl} {\Third.\II}-shout-{\Pl}-{\Prog} {\Obl}-{\DetFour} forest\\
\trans `The men were shouting in the forest.' \hfill (mocCA160713\_2: 01:01:55) \label{ex:determinerplural} 
 \end{exe}
 
 \begin{exe}
\ex 
\gll so tapiɲik i-waɢan-ɡit=oʔ naʔɢa qopaq so l-jaːle-qa   \\
{\DetTwo} armadillo {\Third.\II}-hit-{\LocThree}={\Evd/\Tprl}? day tree {\DetTwo} {\Third.\Poss}-descendent-{\Pl}    \\
	\trans `...the armadillo hit the tree with his children (to get honey)...'
\trans 	\hfill  (mocCA191025: 00:17:21)\label{ex:verbstructure1}
\end{exe}
 
 
 For sake of brevity and relevance to this work, I focus here on the verb complex, whose structure is detailed in \autoref{fig:verbstructure in Northern Chaco Mocoví.}. Special attention is paid to bound person form markers (positions -1 and 3) as they help to draw a major morphological distinction between verb classes.  

%This section deals with basic clause types, showing the typical valence-type predicates and the expression of their core arguments. The semantic and structural criterion for clause categorization will be relevant to identify the P-demotion effect and sets the annotation schema followed in \autoref{sec:annotation-scheema}. 

\begin{figure}[h]
\centering
\scriptsize
\setlength{\tabcolsep}{1pt} % Reduce column padding only inside this group
\begin{tabularx}{\textwidth}{{@{}lllcccccccccccc|cc@{}}}
\\\toprule\toprule
	\textbf{-3}& \textbf{-2} & \textbf{-1} 	&	\textbf{0} &	 \textbf{1} &\textbf{2}	&	\textbf{3}	&\textbf{4} &\textbf{5} & \textbf{6} & \textbf{7} & \textbf{8} &\textbf{9} & \textbf{10} & \textbf{11}& \textbf{12} & \textbf{13}\\\midrule
	
 \textsc{neg}-&\textsc{subj.def}- & \textsc{b.form-} & \textsc{verb} & \textsc{-vmI} & \textsc{-vmII}&\textsc{-b.form}& \textsc{-ben/-trvz} & \textsc{-asp}& \textsc{-des}& \textsc{-refl/-recp}& \textsc{-loc} & \textsc{-dir} & \textsc{trvz} & \textsc{-non.sbj.pl} &\textsc{=dim} & \textsc{=evd/tprl?}\\
 
 \textit{sV-} & \textit{qa-} & \textit{Set I}- & \textsc{root} & \textit{-ɢat}& \textit{-it} & \textit{-Set I}& \textit{-m} & \textit{-ta} & \textit{-ake} & \textit{-taʔ} & \textit{-lek} & \textit{-wek} & \textit{-a(ʔa)}  & \textit{-lo}& \textit{=okiʔ} & \textit{=oʔ} \\
 
 & & \textit{Set II-} &  & \textit{-ɢan} & \textit{-ɢan} & \textit{-Set II} & \textit{-a(ʔa)} & \textit{-sa} & & \textit{-ltaʔ} & \textit{-ɡe} & \textit{-ʃiɡem }&  & & \textit{=oʎiʔ} & \\
 
  & & \textit{Set III-}&  &\textit{-taɢan} &\textit{-taɢan} & \textit{-Set III} & &  \textit{-tak}  & &  &\textit{-ɡi}  &\textit{-o} &  & & \\
 
 & &  &  &  &  &  &  &\textit{-sak} & & & \textit{-ɡit} & \textit{-ɲi}   &  &  &  & \\
 
 & &  &  &  &  &  & & \textit{-teɡ} & &  &  &  & &  &  &
\\\bottomrule\bottomrule
	\end{tabularx}
\bigskip
\caption{Northern Chaco Mocoví verb structure.}\label{fig:verbstructure in Northern Chaco Mocoví.}
\end{figure}


\subsection{Inflectional classes}\label{sec:inflection-classes}
 There are three paradigms of bound person forms \citep[in the sense of][]{Haspelmath2013}, which are referred to in this work as Set I, Set II and Set III -- throughout this work the inflectional class is indicated in the gloss of person forms. As illustrated in (\ref{ex:inflectional classes}), the third person from Set I and Set II includes multiple exponents for this grammatical category and  the distribution of some of them is relevant for detecting the basic valence pattern of predicates (see \ref{sec:argument expression}). The third person also shows partial syncretism and differences that are particular to each paradigm. Note that the form \textit{n-} is shared by Set I and Set III, whereas \textit{i-}, \textit{ɾ-} and \textit{∅-} are identical in Set I and Set II. 
 
\begin{exe}
\ex Sets of bound person forms \label{ex:inflectional classes}
\begin{xlist}
\ex 
 \textbf{Set I}: \textit{dʒ-} `\First\Sg', \textit{qaɾ-} `\First\Pl', \textit{ɾ-...-iɾ} `\Second\Sg', \textit{ɾ-...-iɾ} or \textit{ɾ-...-iː} `\Second\Pl', \textit{i-}, \textit{ɾ-},\textit{∅-}, \textit{n-} `\Third' 

\ex
\textbf{Set II}: \textit{s-} `\First\Sg', \textit{s-...-ɢ} `\First\Pl', \textit{-iɾ} `\Second\Sg', \textit{-iɾ} or \textit{-iː} `\Second\Pl', \textit{i-}, \textit{ɾ-},\textit{∅-}, \textit{t-} `\Third'

\ex
 \textbf{Set III}: \textit{ɲ-} `\First\Sg', \textit{ɲ-...-ɢ} `\First\Pl', \textit{n-...-iɾ} `\Second\Sg', \textit{n-...-i} `\Second\Pl',  \textit{n-} `\Third' 
\end{xlist}
\end{exe}


From a morphological point of view, verbs can be grouped according to the inflectional paradigm they follow and the possibility of alternating between them. The selection of these bound paradigms is lexically determined; thus, each predicate will select a bound person form on a verb-by-verb basis. An internal four-way distinction of types of verbs can be created (for a similar verb classification, see \citet[][96ff]{grondona1998}): (i) predicates that only take Set I, (ii) predicates that only take Set II, (iii) predicates that only take Set III and (iv) predicates that can alternate between Set II and Set III (see \citet{Zurlo2016b} for a detailed study of this phenomenon in Toba and \citet{walker.vanlier2026} for a comparative approach). Here I will not say much about the alternation between paradigms, as it is more relevant for this work to observe the verb type and the verb valence information expressed by the third-person bound form. Observe the contrast between first- and third-person in (\ref{ex:comeback})-(\ref{ex:sit}). 

\begin{exe}
\ex\label{ex:comeback}
Set I with one-place predicate: third-person \textit{ɾ-}
\begin{xlist}
\ex 
	 \gll ajim \textbf{dʒ}-apil    \\
			        {\First\Sg.\Pron} {\First\Sg.\I}-come.back\\
	\trans `I come back (here).' \hfill (mocCA120711: 00:05:10)\label{ex:comeback1}
	
\ex 
	  \gll so-ma jale \textbf{ɾ}-apil \\
			      {\DetTwo}-ma man {\Third.\I.\Intr}-come.back\\
	  \trans `The man comes back (here).' \hfill (mocCA120711: 00:13:52)\label{ex:comeback3}
\end{xlist}
\end{exe}


\begin{exe}
\ex\label{ex:cut}
Set II with two-place predicate: third-person \textit{i-}
\begin{xlist}
\ex 
	 \gll   ajim \textbf{s}-tʃaq-tak so l-aʔat \\
			        {\First\Sg.\Pron} {\First.\II}-cut-{\Prog} {\DetTwo} {\Third.\Poss}-meat\\
	\trans `I'm cutting the meat.' \hfill (mocCA120711: 00:12:54)\label{ex:cutfirst}
	
\ex 
	 \gll so jale \textbf{i}-tʃaq-tak so l-aʔat\\
			      {\DetTwo} man {\Third.\II}-cut-{\Prog} {\DetTwo} {\Third.\Poss}-meat\\
	 \trans `The man is cutting the meat.' \hfill  (mocCA120711: 00:09:24)\label{ex:cutfirst}
\end{xlist}
\end{exe}

\begin{exe}
\ex\label{ex:sit}
Set III with one-place predicate: third-person \textit{n-}
\begin{xlist}
\ex 
	 \gll \textbf{ɲ}-qaɢan-ɲi    \\
			        {\First.\III}-sit-{\Dir:\Down}\\
	\trans `I take a seat.' \hfill  (mocCA120626: 00:36:52) \label{ex:sitfirst}
	
\ex 
	  \gll \textbf{n}-aqaɢan-ɲi \\
			     {\Third.\III}-sit-{\Dir:\Down}  \\
	  \trans `He takes a seat.'
	  \hfill (mocCA120626: 00:39:31)\label{ex:sitthird}
\end{xlist}
\end{exe}


\subsection{Argument expression and clause types} \label{sec:argument expression}
Following a traditional view on clause classification, here I define types of clauses by the valence profile of their predicates  \citep{Dryer2007,Mithun2005,Mithun2005a} and their (morpho)syntactic patterning. This latter feature is in turn recognized by the way in which arguments are treated with respect to verb indexing, word order and oblique flagging. \citep[see][3-4]{comrie.etal2015}. The (morpho-)syntactic pattern of a clause will then serve to recognize clause types such as intransitive and transitive clauses, among other language-specific subtypes. From a typological angle, it is often mentioned that certain classes of predicates are taken to be present in almost all languages, namely mono-valent and bi-valent predicates, whereas  other types of predicates are subject to more language-specific variation, e.g. avalent predicates \citep[][]{Malchukov2010}. Mocoví, and to a large extent its Guaycuruan sister languages, is distinctive in that predicates are lexically categorized as either mono-valent or bi-valent and that these predicate fall into different morphosyntactic patterns. Arguments, in turn, are defined by the the number of semantic participants required by the predicate and their specific morphosyntactic expression to form a grammatical clause \citep{Song2015}. In what follows, we will see that both verb valence and morphosyntactic encoding of participants work together to recognize core arguments.  

 Let us begin with typical one-place predicates. Examples of this type of predicate are \textit{-oʔon} `get married' or \textit{-inta} `be similar' in (\ref{ex:clause_intrans}) and they require a single participant. The main expression of this participant is the bound person form, i.e., \textit{s-}, which encodes the  syntactic function S. Commonly,  co-nominals or pronominal elements of S arguments tend to be overt only for referent identification. Less relevant for this paper are verbless or non-verbal clauses like \textit{jim ɢoɾatom}, where the juxtaposition of elements suffice for a predication.  

\begin{exe}
 
 \ex One-place predicate clause   
  \gll \textbf{s}-\textbf{oʔon} qam qajka ka \textbf{s}-\textbf{inta} jim ɢoɾa-ta-om \\
 {\First.\I}-get\_married but {\Exist.\Neg} {\DetOne} {\First.\II}-be\_similar {\First.\Sg.\Pron} be\_poor-{\Dur}-{\Ints} \\
 \trans ‘... I got married, but I was not the same, I was very poor...' 
 \trans \hfill (mocCA190924\_1: 00:02:20) \label{ex:clause_intrans}
\end{exe}
 


 Two-place predicates  are like (\ref{ex:clause_trans}) and (\ref{ex:two-place_p3rd}), which require two participants, semantically categorized as agent-like and patient-like respectively, and are morphosyntactically expressed in different ways.  The A argument is verb-indexed and the P argument is overtly expressed  by an independent pronoun or nominals. Syntactically, P arguments tend to precede the predicate when they correspond to speech act participants. On the contrary, third-person P arguments follow the verb and their co-nominal or pronominal expressions can occur as in (\ref{ex:two-place_psap}-\ref{ex:two-place_psap_2}).
 
 \begin{exe}
 \ex  Basic two-place predicate clause: 2nd person P
 \gll \textbf{qamiɾ} \textbf{s}-alawat-\textbf{ɢ} \textbf{qamiɾ} \textbf{s}-aʔɡin-\textbf{ɢ} \\
 {\Second\Sg.\Pron} {\First.\II}-kill-{\Pl} {\Second\Sg.\Pron} {\First.\II}-eat-{\Pl} \\
 \trans ‘...we are going to kill you and we are going to eat you...' \hfill(mocCA191025: 00:38:33) \label{ex:clause_trans}
\end{exe}


\begin{exe}
\ex Basic two-place predicate clauses: 3rd person P \label{ex:two-place_p3rd}
\begin{xlist}
\ex
\gll \textbf{s}-oʔkowan-tak \textbf{na} \textbf{isiɡijak} \\
	{\First.\I}-chase-{\Prog} {\DetThree} animal \\
	\trans `... I (usually) was chasing animals...' \hfill{(mocCA191022\_2: 00:03:10)} \label{ex:two-place_psap}
	
	\ex 
	\gll \textbf{s}-aɢan-ɲi \textbf{so} \textbf{i}-\textbf{oɢon}-\textbf{a}  \textbf{n}-\textbf{etʃikse}\\
	{\First.\I}-leave.alone-{\Dir:\Down} {\DetTwo} {\First.\Poss}-take-{\Nmlz.\Obj.\F} {\Indet.\Poss}-lizard \\
	\trans `...I left alone my prey, the lizard...' \hfill{(mocCA191022\_2: 00:03:36)} \label{ex:two-place_psap_2}
\end{xlist}
\end{exe}

In addition to the transitive and intransitive basic clauses, there are other different clause types, defined by with what I call here extended arguments \citep[the terminology is inspired by][3-4]{Dixon2000}. In these cases, both basic one-place and two-place predicates can take an specific type of spatial morphology or other suffixes expressing P arguments. The spatial morphology class, in particular, includes the locative-orientative suffixes, which can express not only topological information but also work similar to what is typologically defined as applicative morphology \citep[for an updated discussion of earlier and more recent descriptions of spatial morphology, see][]{juarez2025}.  Observe the contrast in (\ref{ex:one-place-extended}), where the presence of the locative-orientative \textit{-ɡe} with a one-place predicate allows the expression of a participant that could otherwise be included as an oblique-marked nominal. Predicates like these in (\ref{ex:one-place-extended-b}) deviate from the typical one-place predicates in that the addition of spatial morphology calls for an (otherwise) peripheral participant -- namely, a location that is semantically not required by the verb valence -- as a P argument, whose nominal expression is similar to those in (\ref{ex:two-place_p3rd}) above. Thus, in clauses of this type, the argument extension is noted by the additional verb marking that implies a new P argument and by the non-oblique expression of the implied participant. It should also be noted that predicates of this kind are by default categorized  as one-place predicates, hence the intransitive glossing for the original S argument of the basic clause is maintained.  

\begin{exe}
\ex Extended one-place predicate clause \label{ex:one-place-extended}
\begin{xlist}
\ex
\gll qomiɾ \textbf{qaɾ}-oʔtʃi (ke-na nadʒik) \\
	{\First\Pl.\Pron} {\First.\Pl.\I.\Intr}-be\_afraid {\Obl}-{\DetThree} track \\
	\trans `We are afraid of that track' \hfill{(mocCA120706: 00:50:42)} \label{ex:one-place-extended-a}
	
	\ex
\gll qomiɾ \textbf{qaɾ}-oʔtʃi-\textbf{ɡe} (\Ungr ke)-na nadʒik \\
	{\First\Pl.\Pron} {\First.\Pl.\I.\Intr}-be\_afraid-{\LocThree}	 {\Obl}-{\DetThree} track \\
	\trans `We fear that track' \hfill{(mocCA120706: 00:51:34)} \label{ex:one-place-extended-b}
\end{xlist}
\end{exe}


 % More specifically, the suffix \textit{-ɡe}, for example, entails the locative participant as a core argument of the clause, while keeping the morphological intransitive expression of the S argument. As noted in this example, the extended P argument is expressed morphologically and lexically. Clauses of this type have received a great deal of attention in Guaycuruan languages  and the function of spatial morphology in contexts like these brought novel data to the typologically unusual link between the domain of spatial morphology and valence increasing. 
 
Building on the morphosyntactic evidence presented in (\ref{ex:one-place-extended}), cases like (\ref{ex:oneplace-extended-go}), which includes a movement predicate, also fall within the domain of extended one-place predicates. Additionally, predicates of this type contrast with semantically similar predicates of the kind in (\ref{ex:move-directional}). Unlike predicates with locative-orientative markers, one-place predicates including directional morphology naturally occur with a peripheral locative participant expressed as an oblique. 
   
\begin{exe}
 \ex Extended one-place predicate clause
 \gll t-a-\textbf{ɡe} \textbf{na} \textbf{ʔoːtʃi} ken ʔwe ka-ta n-akit-ɢak n-akit-take l-oq\\
  {\Third.\II.\Intr}-go-{\LocThree} {\DetOne} forest  {\Adv} {\Exist} {\DetFive}-{?} {\Third.\Indet.\Poss}-hope-{\Nmlz} {\Third.\Indet.\Poss}-hope-{\Des} {\Third.\Poss}-food\\
 \trans ‘...she usually went to the forest hoping to receive food...'  
 \trans \hfill (mocCA191010\_1: 00:02:48) \label{ex:oneplace-extended-go}
\end{exe}


\begin{exe}
\ex \label{ex:move-directional}
\gll qamiɾ k-iɾ-\textbf{ʃiɡem} \textbf{ke}-na qopaq \\
 {\Second\Sg.\Pron} move-{\Second.\Sg.\II}-{\Dir:\Up} {\Obl}-{\DetThree} tree 
 \\
 \trans `Did you climb the tree?' \hfill (mocCA120717: 00:36:28)
\end{exe}


As mentioned before, two-place predicates can also combine with locative-orientative morphology or other suffixes to express P arguments. The effect that locative-directional suffixes have on two-place predicates varies, as these suffixes can simply provide spatial information regarding  the P argument or actually extend the number of arguments. Other suffixes expressing a P argument merely serve to indicate the presence of such argument, indexing non-singular number of the P argument. The chain of clauses in (\ref{ex:extended_twoplace_a}) illustrates the use of both P argument expressing morphology. Here the speaker narrates about his past hunting experience and what kind of animals he used to hunt; in this context, the connected clauses should be viewed as statements about general, recurring, former events.  

On the one hand, the locative-orientative suffix \textit{-ɡit} entails that the P argument is coming to the A argument and contributes to the change in meaning of the verb root i.e. \textsc{see} $\rightarrow$ \textsc{find}. Strictly speaking, there  is no addition of a new argument via the spatial marker (as with one-place predicates), but rather a semantic elaboration on how the P argument interacts with the A. On the other hand, the sequence \textit{a-lo} simply expresses the multiplicity of entities in the reference of the P nominal. Although both instances of the nominal `lizard' do not receive any plural marking, the verb morphology indicating non-singular referentiality contribute to stablish the number contrast given in the translation.\footnote{To better understand the English translation offered in (\ref{ex:extended_twoplace_a}), the semantic distinction between \textit{kind-referring} and \textit{object-referring} terms can be particularly helpful. There is a large discussion in the field of semantics on these terms -- which is part of a topic far beyond the scope of this paper --, thus, here I follow ideas in \cite{krifka.etal1995}. The use of \textit{ana ntʃikse} with the predicate \textit{ɲawangit} appears to be of the \textit{kind-referring} expression, as the reference is not to the specific animal, but the genus `lizard'. By contrast, the second use of the same nominal with the predicate \textit{sakontaalo} is \textit{object-referring} in the sense that reference narrows down to the specific animals contained in the \textit{kind}. Furthermor, consider that natural discourse usually shows unconditioned variation in gender and number expression in nominals. Thus, the second use of \textit{na tʃikse} should not be interpreted as referring to a single, indentifiable `lizard', but rather to many of them \citep[see also][]{krifka2025}}  

\begin{exe}
\ex Extended two-place predicates \label{ex:extended_twoplace_a}
\gll s-akon-ta-\textbf{a}-\textbf{lo} \textbf{na} \textbf{ʃilqajk}-\textbf{qa} ɲ-awan-\textbf{ɡit} \textbf{a}-\textbf{na} \textbf{n}-\textbf{tʃikse} s-akon-ta-\textbf{a}-\textbf{lo} \textbf{na} \textbf{tʃikse}  \\
{\First.\II}-take-{\Dur}-{\Arg}-{\Arg.\Pl} {\DetThree} otter-{\Pl} {\First.\III}-see-{\LocThree} {\F}-{\DetThree} {\Indet.\Poss}-lizard {\First.\II}-take-{\Dur}-{\Arg}-{\Arg.\Pl} {\DetThree} lizard\\	
\trans `...(back then when I hunt) I took otters, I found (a) lizard and I took the lizards...'  \trans \hfill(mocCA190924\_1: 00:03:30)
\end{exe}


 Natural discourse examples presented in this section show that while the S and A arguments are invariably indexed in the verb complex, the overt realization of a P argument can be either only nominal or pronominal, or can be co-expressed by verbal morphology and its corresponding nominal or pronominal expression. This argument expression variability is relevant for the data annotation and is implemented for capturing the range of possibilities for expressing P argument in discourse.  
 
\subsection{The expression of P arguments: nominal determiners and their uses}
 A key feature of the nominal expression of core arguments involves the preceding nominal determiner.\footnote{These grammatical elements have been referred to differently in the previous studies on Mocoví and  the other Guaycuruan languages. They are labelled `classifiers' in Pilagá (\citealt[][112-113]{Vidal2001}, \citealt{payne.vidal2020}) and `demonstratives' in Toba \citep[][147-149]{Censabella2002} and Kadiwéo \citep[][61-63]{Sandalo1997}. In Mocoví specifically, the labels of `deictic roots' \citep[][79]{grondona1998}, `classifiers'  \citep[][174-187]{gualdieri1998} and `demonstratives' \citep{carrio2009} have been used to refer to what I call determiners here.} The closed class of six nominal determiners encodes different semantic parameters related to the nominal referent and can be selected on speech context basis. As a whole, the clustering of semantic parameters defining the paradigm of nominal determiners represents a salient typological feature of the nominal domain \citep[cf.][178-181]{Aikhenvald2000}. The semantic parameters relevant for nominal determination are presence vs. absence, position (lying, standing or sitting) and movement/proximity (coming to the speaker or away from the speaker). 
 
 \begin{tabularx}{\textwidth}{cccl} 
\caption{Northern Chaco Mocoví nominal determiners.}\\\toprule\toprule
\multicolumn{2}{c}{\textsc{semantic dimensions}} & \textsc{determiner} 		& \textsc{full gloss}\\\midrule
\textbf{Absence}		&	 													&	 	\textit{ka}	 	&  \DetOne: non-visible, absent\\
\textbf{Presence}		& \textbf{Movement}				 &	\textit{so}		& \DetTwo: going, far\\
										&														& \textit{na}		& \DetThree: coming, close \\
										& \textbf{Position} 					& \textit{ɾa} 		& \DetFour: static, vertical\\
										&				 										&	\textit{dʒi}	& \DetFive: static, horizontal\\
										&				 										&	\textit{ɲi}	& \DetSix: static, sitting \\	
\bottomrule\bottomrule
\label{tab:NCMnominaldeterminers}
\end{tabularx}

Nominal determiners also contribute to the definiteness of nominal referents. In this regards, defined and highly identifiable nouns follow  nominal determiners, as in (\ref{ex:determiner_p_ka}). Furthermore, the selection of determiners does not seem to obey exclusively to the semantic properties of the nominal referent, as discourse factors also appear to play a role. Note in the example below, for example, that the selection of the nominal determiner could have been \textit{ɲi}, as referents of this kind are commonly conceived in a sitting position. Instead, the speaker selected \textit{ka}, which encodes absence of the referent. This nominal determiner fits the discourse context, however, as the speaker narrates events from a remote past, i.e., childhood, which naturally implies that the nominal referent is not present at the moment of speech.   

\begin{exe}
\ex \label{ex:determiner_p_ka}
\gll sa-s-awan-qa-pe-ge-a \textbf{ka} \textbf{malaʔ} \\
{\Neg}-{\First.\II}-see-{\Pl}-{\Prog}-{\LocThree}-{\Obj.\Arg} {\DetOne} bed \\
\trans `...we didn't know the bed...' \hfill(mocCA191122\_2: 00:01:48) 
\end{exe}

However, depending on the referent and, to some extent on the predicate type, nominal determiners can be omitted, which also affects the referential properties of the noun. Example (\ref{ex:no_determiner}) shows that the noun \textit{toːɾta} occurs without a nominal determiner in the eating event. The reference here is not to a single entity, but to multiple entities, which is reinforced by the dinamicity of the eating event in its progressive form. Change in the referentiality of determined and undetermined nouns is more obvious if we compare the expression of A and P arguments. Whereas the former occurs with the determiner \textit{so} and the referent is a single, individualized entity, the latter lacks a determiner and its referentiality is less individualized.
 
\begin{exe}
\ex \label{ex:no_determiner}
\gll \textbf{so} ɾeɢat ∅-keʔe-tak \textbf{toːɾta}{\Sploan}  kemaɾ-aɢa-j\\
{\DetTwo} jaguar {\Third\II}-eat-{\Prog} bread be\_full-{\Nmlz}-{\Attr}-{\F}\\
\trans `...(they say that) the jaguar was eating fried bread and got full...'
\trans \hfill (mocCA191127\_2: 00:09:18)
\end{exe}


Lastly, it is worth mentioning that nominal determiners have a broader range of uses and constitute one of the devices available in the language for discourse linking. Of particular relevance for this work is their clause-combining function, where they  introduce  a noun-modifying clause similar to relative clauses. This property of nominal determiner is not exclusive to Mocoví, as it has also been reported for other Guaycuruan languages (\citealp[114-115]{gualdieri1998}; \citealp[for an overview on Toba relative clauses, see][]{carpio.censabella2012,messineo.porta2009}). In (\ref{ex:det_noun_modifying}), for example, the nominal determiner \textit{so} introduces the extended one-place predicate. This clause modifies the P argument \textit{noʔɢoɾik} `poverty', which in turn expresses a locative role within the modifying clause (see also  (\ref{ex:p-expression-indef}) where a nominal determiner introduces an object complement-like clause).  

\begin{exe}
\ex  \label{ex:det_noun_modifying}
\gll eko s-oleːnta-ɢ-tak n-oʔɢoɾ-ik \textbf{so} s-atʃikqo-ta-\textbf{ɡe}  \\
{\Disc} {\First.\II}-remember-{\Pl}-{\Prog} {\Indet.\Poss}-be\_poor-{\Nmlz.\Obj.\Sg} {\DetTwo} {\First.\III}-descend.{\Pl}-{\Dur}-{\LocThree} \\	
\trans `...we are thinking about the poverty from where we are coming...'
\trans  \hfill(mocCA190924\_1: 00:02:08)
\end{exe}

Examples like (\ref{ex:det_noun_modifying}) illustrate the complexity involved in the expression of P arguments in natural speech.  Whereas  the P argument of the two-place predicate is only expressed lexically, the same argument of the extended one-place predicate is morphologically encoded, referring back to its co-nominal in the main clause. 


%Include here the predicate \textit{talk}. There are near minimal pairs in the discourse where applicative or other elements introduce an argument and previously mentioned.  This discussion is relevant for the type of arguments that can be part of those detransitivized predicates. 
\subsection{Alignment}
Another relevant typological feature of Mocoví is its alignment system. Nouns and pronouns do not distinguish arguments according to their syntactic function (S, A, or P). However, alignment distinctions emerge in argument indexing through the bound person forms described in \ref{sec:inflection-classes}. \citet{juarez2013} provides a detailed analysis of this domain, to which readers are referred for further discussion. The study argues that Mocoví exhibits two major alignment patterns, nominative--accusative and tripartite, distributed according to grammatical person. First- and second-person forms follow a nominative--accusative pattern (\ref{ex:nominative}), whereas third-person forms display tripartite alignment (\ref{ex:tripartite}). 

\begin{exe}
\ex\label{ex:nominative}
First person: nominative-accusative (S=A: \textit{s-} $\neq$ P: no indexed)
\begin{xlist}
\ex 
	 \gll ajim \textbf{s}-asahmat-tak    \\
			        {\First\Sg.\Pron} {\First\II}-cough-{\Prog}\\
	\trans `I'm coughing.' \hfill (mocCA120716: 00:00:07)\label{ex:nominative_a}
	
\ex 
	  \gll  \textbf{s}-etʃaqa-tak so laʔat  \\
			      {\First.\II}-cut-{\Prog} {\DetTwo} meat\\
	  \trans `I'm cutting the meat.' \hfill (mocCA160706: 00:21:06)\label{ex:nominative_b}
	  
\ex 
\gll so-ɢoɾeʔ \textbf{jim} i-waɢan \\
	{\DetTwo}-ɢoɾeʔ {\First\Sg.\Pron} {\Third.\II}-hit \\
\trans `He hit me.' \hfill (mocCA120628: 01:25:04)\label{ex:nominative_c}
\end{xlist}
\end{exe}

\begin{exe}
\ex\label{ex:tripartite}
Third person: tripartite (S: \textit{ɾ-} $\neq$ A: \textit{i-} $\neq$ P: no indexed)
\begin{xlist}
\ex 
	 \gll ɲi-maɢaɾe \textbf{ɾ}-oʔo-tak    \\
	{\DetSix}-{maɢaɾe} {\Third.\II.\Intr}-be/get\_angry-{\Prog}\\
	\trans `He is angry.' \hfill (mocCA120702: 01:50:07)\label{ex:tript_a}
	
\ex 
	  \gll  jim \textbf{i}-otawan-tak   \\
			 {\First\Sg.\Pron}     {\Third.\II}-help-{\Prog} \\
	  \trans `He is helping me.' \hfill (mocCA180719: 00:19:09)\label{ex:tript_b}
	  \newpage
\ex 
\gll ajim s-elaɾ \textbf{so} \textbf{i}-\textbf{jaːle}-\textbf{k}  t-aw-ɡe ɾa pweblo{\Sploan} \\
	{\First\Sg.\Pron}  {\First.\II}-send/order {\DetTwo} {\First\Sg.\Poss}-descendent-{\M}  {\Third.\II.\Intr}-go-{\LocThree} {\DetFour} town \\
\trans `I sent my child to town.' (Lit. `I sent my child, he goes to town.')
\trans  \hfill (mocCA160722\_1: 01:11:06)\label{ex:tript_c}
\end{xlist}
\end{exe}

\section{Agent foregrounding in the context of valence alternations}
After reviewing the morphosyntactic properties of argument expression, we now turn to the means by which P arguments can be demoted. We focus here on derivational elements \textit{-ɢan} and \textit{-taɢan}, which occur in positions 1 and 2 in \autoref{fig:verbstructure in Northern Chaco Mocoví.}. Such derivational elements can both increase and decrease verb valence, creating constructions that fall within the domains of causativization and antipassivization. These dedicated valence modifiers primarily trigger changes in the morphosyntactic expression of S and A  arguments. That is, the S argument functions as a P argument when one-place predicates increase their valence and the A argument becomes S when two-place predicates reduce their valence. Consequently, these changes in argument structure naturally affect the semantic properties and morphosyntactic expression of P arguments.  The typologically unusual co-expression pattern has been studied by scholars working on Mocoví \citep{Carrio2015,Carrio2015a,Juarez2017c,juarez.alvarez-gonzalez2021} and is attested in other Guaycuruan languages as well. The relevance of this phenomenon relies in that only a  handful of languages have been reported to display this specific valence co-expression pattern. In a recent typological study on voice syncretism, for example,  \citet[][100-103]{bahrt2021} listed Soninke (Western Mande), Makalero (Timo Alor Palantar), Mocoví (Guaycuruan), Wolof (Atlantic Congo)  and Ganja Balanta (Atlantic Congo), as the languages showing what the author classifies as causative-antipassive syncretism.\footnote{However, it should be mentioned that this syncretism is also present in other Guaycuruan languages, e.g., Toba, including its different geographical variants, namely Eastern and Western Formos Toba, Paraguay Toba, as well as the extinct Abipón and the Mataguayan language Nivaclé.}  

When \textit{-ɢan} and \textit{-taɢan} function as intransitivizing devices, the resulting derived predicates are part of what I call agent foregrounding constructions. I employ this language-specific label as it coherently captures the common semantic and structural denominators in the dual function of these verbal markers. %More importantly, it emphasis the fact that is the agent-like clause and its action what these clauses are about. 
    From a more systemic view, the agent foregrounding label makes explicit the functional and structural contrast with other construction types such as the subject defocusing construction, whose dedicated marker is \textit{qa-} and occurs in position -2 (see below (\ref{ex:defckill})).  Inspired by the framework on passives and related constructions in \citet{Shibatani1985a}, I use the notions of foregrounding and defocusing  as functional umbrella terms, indicating that subject and agent-like arguments either increase or decrease their discourse prominence. Functionally,  foregrounding and defocusing constructions mirror each other with respect to the effect that they have on the subject argument of these constructions.  

In agent foregrounding constructions, specifically, the agent-like argument becomes structurally and discursively highly prominent. A by-product of such a change in argument prominence is the demotion of the P argument, which structurally manifests in different ways, as we will see below (see  \ref{sec:structural_feautres_p_dem}). An example of an agent foregrounding construction is given in (\ref{ex:objsupkill}). In this example, the first use of the predicate \textit{kill} inflects intransitively and the entailed P argument is syntactically omitted. The predication is only about the agent's activity and does not specify a clear, bounded outcome.  

\begin{exe}
\ex \textit{-ɢan} agent foregrounding construction
	 \gll so pjoq i-taː-taʔ-pe-ɡe so ʎik so jale \textbf{ɾ}-lawat-\textbf{ɢan} \textbf{i}-lawat a-\textbf{so} \textbf{una}{\Sploan} \textbf{neketeɡeɢaj}   \\
	{\DetTwo} dog {\Third.\II}-sniff-{\Pl}-{\Prog}-{\LocTwo} {\DetTwo} track {\DetTwo} man {\Third.\II.\Intr}-kill-{\Vm:\Intr}   {\Third.\II}-kill {\F}-{\DetTwo} {\Det} girl   \\
	\trans `The dog was sniffing the track of \textbf{the killer} (who) killed the girl.'
	\trans Lit. `The dog was sniffing the track, the man kills, he killed a girl.'
		\trans\hfill (mocCA130614\_1: 00:09:42)\label{ex:objsupkill}
\end{exe}

Furthermore, the agent foregrounding construction is both structurally and semantically connected to causative constructions, as the contrast in  (\ref{ex:causefall}) illustrates. Morphologically, only a small class of one-place predicates can be causativized by \textit{-ɢan}, which introduces a cause participant, functioning as the A argument. This causativization process typically requires the bound person form \textit{i-}, which indexes third-person, transitive subjects. Semantically, both the agent foregrounding and the causative constructions involve an agent-like argument that functions either as S (in the detransitivized clause) or A (in the causative clause) \citep[see][Ch.4]{juarez2023}   

\begin{exe}
\ex  \textit{-ɢan} causative construction \label{ex:causefall}
\begin{xlist}
\ex
\gll n-wiɾ so waɢajaɢ jaɢat ∅-ʔɡen-ɲi so iːmek \\
{\Third.\III}-come {\DetTwo} water rain {\Third.\II}-fall-{\Dir:\Down} {\DetTwo} house\\
\trans `The water, the rain, came, and the house fell down.'
 \hfill  (mocCA160720: 00:00:49)\label{ex:causefalla}


\ex
\gll so jaɢat \textbf{i}-ʔɡen-\textbf{ɢan}-ɲi ɲi iːmek\\
{\DetTwo} rain  {\Third.\II}-fall-{\Vm:\Tr}-{\Dir:\Down} {\DetSix} house\\
\trans `The rain knocked down the house.'
\hfill{ (mocCA160720: 00:03:39)} \label{ex:causefallb}

\end{xlist}
\end{exe}

On the other hand, agent foregrounding constructions functionally contrast with the subject-defocusing construction, illustrated in (\ref{ex:defckill}). In this construction, the subject is expressed as a non-referential, impersonal-like participant (e.g., `someone', `they', `people') or as a  plural participant,  and only the P argument is discursively prominent. This clause does not inflect intransitively, as opposed to the agent foregrounding construction,  and its reading(s) is context dependent, always entailing a non-specific subject argument.

\begin{exe}
\ex Subject defocusing
	 \gll dʲoː dʒi i-wa \textbf{qa}-i-lawat \textbf{qa}-i-lawat dʒi i-wa \\
			       {\Exclam} {\DetFive} {\First\Sg.\Poss}-wife {\Sbj.\Defc}-{\Third.\II}-kill {\Sbj.\Defc}-{\Third.\II}-kill {\DetFive} {\First\Sg.\Poss}-wife\\
	\trans `...oh, someone killed \textbf{my wife}, someone  killed \textbf{my wife}...' \trans\hfill (mocCA191127\_2: 00:13:32)\label{ex:defckill}
\end{exe}

 
\subsection{Structural correlates and effects on the P argument}\label{sec:structural_feautres_p_dem}
To understand the morphological and syntactic nuances of agent foregrounding constructions, we begin by examining examples from controlled, elicited sentences. One of the salient features of Mocoví is its person-based alignment pattern, combining nominative-accusative and tripartite alignments. Thus, the structural correlates of the agent foregrounding constructions are slightly different whether the foregrounded A argument corresponds to a speech act participant or not. 

 First- and second-person arguments follow a nominative-accusative alignment at the indexing level -- consider, for instance, A=S: \textit{s-} $\neq$ P: \textit{∅} in (\ref{ex:kill_firstperson}). Because of this alignment type, the foregrounding effect on the subject argument does not correlate with any formal change in its verbal expression. Foregrounding the A/S activity only leads to  disregarding the nominal determiner, as in (\ref{ex:kill_firstperson_b}), therefore focusing on the verb action itself rather than on its outcome. Consequently, the meaning of the verb also differs slightly from the base, unmarked transitive construction. The entailment in (\ref{ex:kill_firstperson_b}), is that the agent has not killed  a pig yet, but he is about to go somewhere else to engage in that specific activity, without knowing whether the number of killed animals will be one or several.  


\begin{exe}
\ex \label{ex:kill_firstperson}
\begin{xlist}
\ex
\gll ajim \textbf{s}-alawat so kos \\
	{\First\Sg.\Pron} {\First.\II}-kill {\DetTwo} pig \\
	\trans `I killed the pig.' \hfill(mocCA160725: 00:35:22) \label{ex:kill_firstperson_a}
\ex
 \gll ajim \textbf{s}-alawat-\textbf{ɢan}  kos \\
	{\First\Sg.\Pron} {\First.\II}-kill-{\Vm:\Intr}  pig \\
	\trans `I'm going to slaughter (a) pig(s).' \label{ex:kill_firstperson_b} \hfill(mocCA160725: 00:35:22)

 \end{xlist}
\end{exe}

As mentioned above, when two‑place predicates are marked by \textit{-ɢan} and the A argument corresponds to a non-speech act participant, these predicates usually call for the bound person \textit{ɾ-}, which indicates an intransitive subject. In addition, the P argument can be fully omitted (\ref{ex:intrtieupb}) or expressed as a bare noun without a nominal determiner (\ref{ex:intrtieupc}).  In both cases, the demoted P argument is interpreted as generic, lacking a sense of uniqueness, although subtle semantic differences remain. In these examples, the complete omission of the P argument allows for a much wider set of entities that could  potentially be tied up, e.g., horses, cows, sheep, and the like, and only contextual clues can help identify the specific entities involved in this event. By contrast, the use of a bare noun narrows the possible referents to a specific class, in this case \textit{horses}.

%All of the examples so far considered show that valence reduction of two-place predicates correlates with an intransitive syntactic structure in that only the subject of the derived predicate can be expressed in the clause. However, there is a syntactic strategy by which the expression of the object argument is allowed. This strategy has to do specifically with the absence of a nominal determiner.

\begin{exe}
\ex\label{ex:intrtieup}
\begin{xlist}
\ex 
	 \gll so jale n-qoin so ʃipeɢaɢ\\
			        {\DetTwo} man {\Third.\III}-tie.up {\DetTwo} horse\\
	\trans `The man tied up a horse.'
	 \hfill (mocCA160725\_1: 00:20:47) \label{ex:intrtieupa}

\ex 
	 \gll so jale \textbf{ɾ}-oqoin-\textbf{ɢan}  \\
 {\DetTwo} man {\Third.\II.\Intr}-tie.up-{\Vm:\Intr}\\
	  \trans `The man goes to tie up (horses).'
	  \hfill (mocCA160725\_1: 00:28:09) \label{ex:intrtieupb}

\ex 
	  \gll \textbf{ɾ}-oqoin-\textbf{ɢan} \textbf{ʃipeaɢ} \\
			  {\Third.\II.\Intr}-tie.up-{\Vm:\Intr}  horse   \\
	  \trans `He goes to tie up horses.'
	   \hfill (mocCA160725\_1: 00:30:25) 	\label{ex:intrtieupc}
	\end{xlist}
\end{exe}

 In \textit{ɢan}-marked two-place predicates, the expression of P arguments preceded by a nominal determiner results in ungrammatical clauses, regardless of whether the referential noun is singular (\ref{ex:intrtieupungra}) or non-singular (\ref{ex:intrtieupungrb}), indicated by plural and collective morphology.\footnote{One would expect a clause like (\ref{ex:intrtieupungrb}) to  be acceptable without the nominal determiner, but I could not test this specific context through elicitation. It should also be mentioned that natural discourse examples in my corpus do not include examples like (\ref{ex:intrtieupungr}) either.}   

\begin{exe}
\ex\label{ex:intrtieupungr}
\begin{xlist}
\ex [\Ungr]{	
	  \gll so jale ɾ-qoin-ɢan so ʃipeaq \\
			 {\DetTwo} man {\Third.\III}-tie.up-{\Vm:\Intr}  {\DetTwo} horse      \\
	 \trans `The man goes to tie up the horse.'	 
	 \trans \hfill (mocCA160725\_1: 00:29:01, rejected clause)}	\label{ex:intrtieupungra}	
	
\ex [\Ungr]{	
	  \gll  ɾ-qoin-ɢan so ʃipeaɢ-ɾ-pi  \\
			    {\Third.\III}-tie.up-{\Vm:\Intr}  {\DetTwo} horse-{\Pl}-{\Coll}   \\
	  \trans `He goes to tie up the horses.'
	  \trans \hfill(mocCA160725\_1: 00:31:26, rejected clause)}\label{ex:intrtieupungrb}	
\end{xlist}
\end{exe}

Since this detransitivized construction requires a non-definite P argument, it naturally follows that its syntactic expression is ruled out when it refers to an individual, uniquely identifiable entity. Compare the derivation of the predicates (\ref{ex:intrhunger}) and (\ref{ex:intrhungerc}). The causativized predicate  allows the expression of the first-person singular P argument. However, when the clause is intransitivized by \textit{-ɢan}, the construction becomes ungrammatical if the independent pronoun is expressed. Furthermore, note that the verb expression of the causee argument is also removed and the valence modifier occupies its slot. 

\begin{exe}
\ex \label{ex:intrhunger}
\begin{xlist}
\ex 
\gll so doqolek siempre{\Sploan} \textbf{jim} \textbf{i}-qopat-\textbf{ɢat-it} \\
{\DetTwo} non.indigenous(person) always {\First\Sg.\Pron} {\Third.\II}-be.hungry-{\Caus}-{\Causee}\\
\trans `The non-indigenous always makes me starve.'
\trans  \hfill (mocCA191205: 00:13:39) \label{ex:intrhungera}
	
\ex 
\gll so doqolek siempre{\Sploan} \textbf{ɾ}-qopat-\textbf{ɢat-ɢan} \\
{\DetTwo} non.indigenous(person) always  {\Third.\Intr.\II}-be.hungry-{\Caus}-{\Vm:\Intr}\\
\trans `The non-indigenous (person) always makes (someone/people) starve.'
\trans  \hfill (mocCA191205: 00:11:59)\label{ex:intrhungerb}
\end{xlist}
\end{exe}

\begin{exe}
\ex [\Ungr] {
\gll so doqolek siempre{\Sploan} \textbf{ajim} \textbf{ɾ}-qopat-\textbf{ɢat-ɢan} \\
{\DetTwo} non.indigenous always {\First\Sg.\Pron} {\Third.\Intr.\II}-be.hungry-{\Caus}-{\Vm:\Intr} \\
\trans `The non-indigenous (person) always makes (me) starve.'
\trans \hfill(mocCA191205: 00:12:40, rejected clause)}\label{ex:intrhungerc}
\end{exe}

\subsection{How to recover the expression of P}
Although agent foregrounding construction indirectly involve P demotion, the language displays several morphosyntactic strategies to overtly reintroduce the demoted argument and narrow its referentiality in the immediate discourse context. These syntactic strategies respond to the functional need of avoiding ambiguity in participant referentiality as the discourse unfolds, specifically regarding the erstwhile demoted P argument.


One of the strategies to encode the P argument of a detransitivized predicate is by the juxtaposition of two predicates that otherwise  function as individual verbs on their own. This combination of predicates produces structures similar in many regards to serial verb constructions \citep[see][]{Aikhenvald2006, Bisang2009, Haspelmath2016}, although I do not treat them as a single complex clause here.  The multiverbal strategy to encode the P argument of a detransitivized predicate can be further divided according to the valence of the combined predicates. On the one hand, we find combined predicates like \textit{ɾ-kin-ɢan} `She/He greets (somebody)' and \textit{i-kin} `She/He greets him/her' in (\ref{ex:objincl1}). This sequence of predicates involves then the \textit{ɢan-}intransitivized predicate combined with a basic two-place predicate, creating a clause with the following transitivity pattern: \textsc{intr} + \textsc{tr}.  
           
\begin{exe}
\ex 
	 \gll ∅-k-pe-ɡe-eʔ ɾ-kin-\textbf{ɢan} i-\textbf{kin} \textbf{dʒi} \textbf{l}-\textbf{jaːle}-\textbf{qa}-\textbf{pi}   \\
			        {\Third.\II}-move-pe-{\LocTwo}-{\Pl} {\Third.\II.\Intr}-greet-{\Vm:\Intr} {\Third.\II}-greet {\DetFive} {\Third.\Poss}-descendent-{\Pl}-{\Coll}\\
	\trans `S/he goes to greet his/her children.'
\trans	 Lit. S/He goes, (S/he) greets, (S/he) greets his/her children.
	\trans \hfill  (mocCA160725: 00:55:30)\label{ex:objincl1}
\end{exe}

 Note that the selection of the second predicate does not appear to be arbitrary, as it repeats the  same verb root, \textit{-kin} `greet', on which the \textit{ɢan-}intransitivized predicates is built. Thus, one could suggest that the greeting event depicted in the translation is made of two predications semantically connected by the shared verb meaning. While the first predication depicts the agentivity of the A argument, the second one helps to introduce the actual referent of the previously demoted P argument. 

Another possibility to express the syntactically omitted P argument is illustrated in (\ref{ex:objincl2}). Unlike (\ref{ex:objincl1}), the first detransitivized predicate \textit{-kin-ɢan} is further derived by the stacking of \textit{-ɢan} in its transitivizing function, which adds a new subject argument to the verb valency. This creates a morphologically complex predicate that combines with another non-derived two-place predicate. Both conjoined predicates show the following transitivity pattern: \textsc{tr} + \textsc{tr}. Like (\ref{ex:objincl1}), the main function of the second predicate is the syntactic expression of the P argument that was part of the verb root employed in the double-derived predicate.  

\begin{exe}
\ex 
	 \gll ajim s-kin-\textbf{ɢan}-\textbf{ɢan} so i-jaːle-k i-\textbf{kin} \textbf{so} \textbf{l}-\textbf{piɾ}   \\
			        {\First\Sg.\Pron} {\First.II}-greet-{\Vm:\Intr}-{\Vm:\Tr} {\DetTwo} {\First\Sg.\Poss}-descendent-{\M} {\Third.\II}-greet {\DetTwo} {\Third.\Poss}-grandfather\\
	\trans `I made my son greet his grandfather.'  \hfill(mocCA160725: 00:46:37)\label{ex:objincl2}
\end{exe}

 The set of examples above includes a derived verb root categorized as a two-place predicate in its basic non-derived use. The multiverb strategy to encode the omitted P argument is also possible if the verb root is of one-place. The main condition, however, is that the combined predicates adjust to one of the transitivity patterns already mentioned, specifically:  \textsc{intr} + \textsc{tr}. To achieve this syntactic pattern, the basic one-place predicate must be further derived.  First,  note that the \textit{taɢan-}intransitivized predicate does not accept the syntactic expression of P, as in (\ref{ex:showernoobj}). But, in order to include it, the addition of the transitivized predicate \textit{-tʃil-ɢan} `shower him/her' is needed (\ref{ex:showerwithobject}). 

\begin{exe}
\ex\label{ex:}
\begin{xlist}
\ex 	
\gll ɾ-tʃil-ɢan-taɢan (\Ungr so qaɾ-qaja-ɾ-pi)\\
     {\Third.\III}-shower-{\Vm:\Tr}-{\Vm:\Intr} {}{\DetTwo} {\First\Pl.\Poss}-brother-{\Pl}-{\Coll}  \\
\trans `S/He baptizes (someone).'  \hfill  (mocCA160726: 00:56:25)\label{ex:showernoobj}
	
\ex 
\gll \textbf{ɾ}-atʃel-ɢan-\textbf{taɢan}-tak n-atʃil-ɢan-tak so qaɾ-qaja-ɾ-pi i-noɢon-o ke-ɾa n-qaːt-ek  \\
{\Third.\III.\Intr}-shower-{\Vm:\Tr}-{\Vm:\Intr}-{\Prog} {\Third.\III}-shower-{\Vm:\Tr}-{\Prog} {\DetTwo} {\First\Pl.\Poss}-brother-{\Pl}-{\Coll} {\Third.\II}-enter-{\Dir:\In} {\Obl}-{\DetFour} {\Indet.\Poss}-speak-{\Sg.\Obj.\Nmlz}\\
\trans `He is baptizing our brothers who joined the Word (of God).'
\trans Lit. `He is baptizing, (he) is baptizing our brothers, (they) went into the Word.'
\trans \hfill  (mocCA160726: 00:46:37) \label{ex:showerwithobject}	
\end{xlist}
\end{exe}

In the combined predicates like (\ref{ex:objincl1}), (\ref{ex:objincl2}), and (\ref{ex:showerwithobject}) the second predicate introducing the P argument appears to be lexically restricted to the same verb root that serves as the base for the \textit{ɢan-} or \textit{taɢan-} detransitivized predicates. 

Finally, another way to encode the former P argument of the detransitivized predicates is by marking it as an oblique, as illustrated in (\ref{ex:objincl3}). This strategy is possible once the verb valence of the base predicate is re-augmented by \textit{-ɢan}, which results in the double-derived verb \textit{-kin-ɢan-ɢan} `made to greet'. Rather than resulting in a clause with two non-subject arguments, this type of double derived predicate includes only two core arguments, the A and the P, in addition to the oblique-marked participant. 

\begin{exe}
\ex 
	 \gll so i-taʔa ajim i-kin-ɢan-ɢan \textbf{ke}-so qaɾ-piɾ     \\
			 {\DetTwo} {\First\Sg.\Poss}-father {\First\Sg.\Pron} {\Third\II}-greet-{\Vm:\Intr}-{\Vm:\Tr} {\Obl}-{\DetTwo} {\First\Pl.\Poss}-grandfather\\
	\trans `My dad made me greet our grandfather.' \hfill (mocCA160725: 01:01:22) \label{ex:objincl3}
\end{exe}
 
It must be mentioned, however, that examples like (\ref{ex:objincl3}), are rather infrequent in natural discourse, but naturally accepted when constructed via elicitation techniques. A much larger corpus-based study should elucidate the discourse factors playing a role in the selection of either the oblique-marked expression of the demoted P argument or the multi-verbal strategy, with the demoted P as a core argument of the repeated verb root. This task is beyond the scope of this paper. 

\section{Lexical distribution and importance for lexicon building}\label{sec:lexicon_distribution}
Since agent foregrounding constructions overlap in many respects with properties associated with antipassive constructions, I draw on recent typological observations on antipassives to account for the distributional properties of the agent foregrounding morphological markers. \autoref{tab:predicates taking -ɢan and -taɢan} summarizes the distribution of the two‑place predicates that occur with \textit{‑ɢan} and \textit{‑taɢan}.

First, note that these elements exhibit complementary and asymmetric distributional patterns. Specifically, predicates that are detransitivized by \textit{‑ɢan} do not occur with \textit{‑taɢan}, and vice versa. Additionally, the shorter form shows a much wider distribution than the longer form, which is restricted to only a handful of predicates. Furthermore, although previous studies have not explicitly addressed the frequency of these suffixes, examples involving \textit{‑ɢan} outnumber those with \textit{‑taɢan} in the available linguistic descriptions of Mocoví \citep[see][]{gualdieri1998,Carrio2015,Carrio2015a}. Such an asymmetric distributional pattern is not unusual in languages that possess multiple antipassive marker \citep[see, for instance,][]{say2021,janic.witzlack-makarevich2021}. \citet[][]{serzant.etal2021}, furthermore, argue that frequency plays a role in the asymmetric distribution of antipassive markers coexisting in the same language. Specifically, they propose that ``the most frequent lexical input will correlate with the less marking on the verb and viceversa".   
%\newpage
%\begin{onehalfspace}
\begin{center}
\small
\begin{tabularx}{\textwidth}{XXXX}
\caption{Two-place predicates taking \textit{-ɢan} and \textit{-taɢan} `\Vm:\Intr'.} \label{tab:predicates taking -ɢan and -taɢan} \\
\toprule\toprule
\multicolumn{2}{l}{\textbf{Two-place predicate}} & \multicolumn{2}{l}{		 								 	 \textbf{\textit{ɢan-}derived predicate}}\\\toprule
\endfirsthead 
\multicolumn{4}{c}{\tablename\ \thetable\ -- \textit{Continued from previous page}} \medskip \\
\toprule 
\multicolumn{2}{l}{\textbf{Two-place predicates}} & \multicolumn{2}{l}{		 								 	 \textbf{\textit{ɢan-}derived predicate}}\\\toprule
\endhead\midrule
\multicolumn{2}{l}{} & \multicolumn{2}{r}{\textit{Continued on next page}} \\
\endfoot
\bottomrule\bottomrule
\endlastfoot
 \textit{i-waɢan} & `S/he/it hits him/her/it' & \textit{ɾ-waɢan-\textbf{ɢan}} & `S/he/it hits at (sth)'\\ 
\textit{i-tʃaq} & `S/he/it cuts him/her/it' & \textit{ɾ-tʃaq-\textbf{ɢan}} & `S/he goes to cut (sth)'\\
\textit{i-kin} & `S/he greets him/her'  & \textit{ɾ-kin-\textbf{ɢan}} & `S/he/it greets (sb)'\\ 
 \textit{i-mit}& `S/he/it searches for him/her/it' & \textit{ɾ-mit-\textbf{ɢan}}& `S/he/it searches around'\\
 \textit{i-lawat} & `S/he/it kills him/her/it' & \textit{ɾ-lawat-\textbf{ɢan}} & `S/he goes to kill'\\ 
  %\textit{∅-ʔɢaja} &`S/he/it listen' &\textit{n-ʔɢaja-\textbf{ɢan}}& `S/he/it obeys' (\textsc{bck}ː)\\
 \textit{i-taː} & `S/he/it smells'  & \textit{ɾ-taː-\textbf{ɢan}} & `S/he/it sniffs (around)' \\
 \textit{i-ɾa} & `S/he writes sth.' & \textit{ɾ-ɾa-\textbf{ɢan}} & `S/he/it writes (sth)' \\ 
\textit{i-ʔɢat} & `S/he talks to him/her.' & \textit{ɾ-ʔqat-\textbf{ɢan}} & `S/he talks' \\ 
\textit{i-meːt} & `S/he  piles sth. up' & \textit{ɾ-meːt-\textbf{ɢan}} & `S/he piles up (sth)' \\ 
\textit{i-laɾ} & `S/he  sends him/her' & \textit{ɾ-lat-\textbf{ɢan}} & `S/he gives orders' \\ 
\textit{i-ʔleːnte} & `S/he  remembers him/her/it' & \textit{ɾ-ʔleːnte-\textbf{ɢan}} & `S/he remembers (sth)' \\ 
\textit{i-lat} & `S/he leaves him/her/it. & \textit{ɾ-lat-\textbf{ɢan}} & `S/he leaves (sth/sb)' \\ 
\textit{i-ʃila} & `S/he asks him/her (sth)' & \textit{ɾ-ʃila-\textbf{ɢan}} & `S/he asks (sth)' \\
\textit{i-nat} & `S/he asks for sth.' & \textit{ɾ-nat-\textbf{ɢan}} & `S/he implores (him/her)' \\ 
\textit{i-wale} & `S/he borrows sth' & \textit{ɾ-wale-\textbf{ɢan}} & `S/he borrows (sth)' \\
\textit{i-qawalat} & `S/he  kneads sth.' & \textit{ɾ-qawalat-\textbf{ɢan}} & `S/he kneads (sth)' \\ 
\textit{n-qat} & `S/he  pulls sth. up' & \textit{ɾ-qat-\textbf{ɢan}} & `S/he goes to harvest' \\ 
\textit{n-maɾen} & `S/he  scrapes sth' & \textit{ɾ-maɾen-\textbf{ɢan}} & `S/he goes to scrape (sth)' \\ 
\textit{n-owaɾ} & `S/he sews sth' & \textit{ɾ-owaɾ-\textbf{ɢan}} & `S/he sews (sth)' \\ 
\textit{∅-peloɡ} & `S/he rakes sth.' & \textit{ɾ-peloɡ-\textbf{ɢan}} & `S/he rakes (sth)' \\ 
\textit{∅-ʔleg} & `S/he sweeps sth.' & \textit{ɾ-ʔleɡ-\textbf{ɢan}} & `S/he sweeps (sth)' \\ 
\textit{∅-kijo} & `S/he washes sth.' & \textit{∅-kijo-\textbf{ɢan}} & `S/he washes (sth)' \\ 
\textit{∅-lapon} & `S/he piles sth.' & \textit{∅-lapon-\textbf{ɢan}} & `S/he piles (sth)' \\ 

 \multicolumn{2}{l}{} & \multicolumn{2}{l}{ \textbf{\textit{taɢan}-derived predicate}} \\\cmidrule{3-4}
 \textit{i-paɢaɡin} & `S/he teaches him/her' & \textit{ɾ-paɢaɡin-\textbf{taɢan}} & `S/he teaches (sth)' \\
%\textit{•} & & \textit{i-keʔe-\textbf{ɢan}} & `S/he/it feeds him/er/it' \\
\textit{i-ʔden} & `S/he knows sth.' & \textit{ɾ-ʔden-\textbf{taɢan}} & `S/he knows (sth)'\\
\textit{n-taɾen} & `S/he cures him/her' & \textit{n-taɾen-\textbf{taɢan}} & `S/he cures (sb)'\\
\textit{n-oʔwen} & `S/he works on sth.?' & \textit{ɾ-oʔwen-\textbf{taɢan}} & `S/he works'\\
\end{tabularx}
\end{center}

%Second, arguments on verb types. Examples of ingestive verbs. 

%Third, there is no conditioning on the type of P argument that is demoted. 

%\end{onehalfspace}

%\section{Foregrounding constructions and uses}
%
%%%%% Copy the structure of examples. They should be in this format. 
% \begin{exe}
% \ex \label{ex: }
% \begin{xlist}
% \ex 
% \gll  \\
%  \\
% \trans ‘.' \label{ex:}
%
%\ex 
%\gll \\
%\\
%\trans ‘.'  \label{ex:}
% \end{xlist}
%\end{exe}

  Second, as far as I am aware, the only  lexical constraint conditioning which predicates can occur in agent foregrounding constructions relates to their basic verb valence. From the verb list above and natural discourse data, it becomes clear that the  relevant parameter is that the base predicate must be a two-place predicate. In fact, the use of \textit{-ɢan} and \textit{-taɢan} as detransitivizers is a natural language-specific morphological test to identify the class of two-place predicates. Regardless of the basic verb valence, it is worth noticing nevertheless that the predicates in \autoref{tab:predicates taking -ɢan and -taɢan} align with some of the properties that \citet[][182]{say2021} describes as characteristics of natural antipassives, including the types of verbs that typically undergo antipassivization \citep[see, the verb hierarchy proposed by][]{malchukov2015}. Semantically, predicates in their basic form cover a large range of verb classes, including not not only proper activities (`greet'), but also contact (`hit'), pursuit (`search'), cognition (`know', `remember') and speaking verbs (`ask', `ask for'), among others. As noted earlier, this verb list nicely illustrates that the base predicate typically includes a highly agentive A argument and the derived predicate corresponds to an activity that highlights the agent's involvement in the event.
  
  However, there are only a few two‑place predicates that deviate from the general derivational pattern associated with detransitivizing functions in agent foregrounding constructions. In particular, two‑place predicates belonging to the class of ingestive verbs -- such as ‘eat’, ‘drink’, ‘smoke’, and ‘lick’ -- only combine with  \textit{-ɢan} as a causativizer rather than as a marker of the agent foregrounding construction. Recall that Mocoví is categorized as transitivizing language which strongly favors causativization of one-place predicates (see \ref{sec:argument expression}).
  
  \begin{exe}
\ex\label{ex:causeeat}
\begin{xlist}
\ex 
 \gll so noɡot=okiʔ ∅-keʔe-tak pan{\Sploan}   \\
			        {\DetTwo} adolescent={\Dim.\M} {\Third.\II}-eat$_1$-{\Prog} bread\\
	\trans `The kid is eating bread.'
	 \hfill (mocCA160726: 00:03:02)

	
\ex 
 \gll ∅-keʔe-\textbf{ɢan} pan{\Sploan} so noɡot=okiʔ\\
 {\Third.\II}-eat$_1$-{\Vm:\Tr} bread {\DetTwo} adolescent={\Dim.\M}\\
\trans `(The woman) fed bread to the kid.'
	  \hfill  (mocCA160726: 00:33:17) 
\end{xlist}
\end{exe}
  
  \begin{exe}
\ex \label{ex:causedrink}
\begin{xlist}
\ex 
 \gll so jale maʃi i-me-wek n-eʔet-tape-ɡi na l-ataɾ   \\
			       {\DetTwo} man already {\Third.\II}-end-{\Dir:\Out} {\Third.\III}-drink-{\Prog.\Pl}-{\LocFour} {\DetThree}  {\Third.\Poss}-medicine\\
	\trans `The man already got better (because) he is taking his medicine.' 
	\trans \hfill(mocCA160711\_1: 00:43:17)\label{ex:causedrinka}

\ex 
	 \gll n-eʔet-\textbf{ɢan} \\
{\Third.\III}-drink-{\Vm:\Tr}\\
	  \trans `S/he/it makes him/er drink it.'
	  \hfill \citep[][115]{Buckwalter1995}\label{ex:causedrinkb}
\end{xlist}
\end{exe}  
  
  
   Typologically, it is not unusual for ingestive verbs to show idiosyncratic behavior with respect to the valence alternation patterns they participate in and the transitivity profile of the verb lexicon \citep[see][]{Nichols2004}. In a typological survey on causative constructions \citet[][63-65]{Dixon2000}  concluded that ``drinking and eating are the transitive activities which people are most likely to make other people do, on every other continent". Similarly, for languages with an intransitive constraint on morphological causativization, \citet[][30-31]{Naess2009} and \citet[][48-54]{Amberber2009} observed more recently that not only are `eat' and `drink' flexible in terms of causativization, but so are ``other verbs referring to acts of  ingestion or consumption'', e.g.  `smoke' and `lick', as well as verbs of perception, e.g. `know'.  
   
   Lastly, given that the main focus of agent foregrounding constructions is the agent’s activity, it naturally follows that the animacy of the demoted P argument is not a relevant parameter of variation. Typological work has shown that the animacy of demoted P arguments plays a role in the variation observed in antipassives  \citep{say2021,zuniga.kittila2019}. However, the list of predicates in \autoref{tab:predicates taking -ɢan and -taɢan} and the examples discussed so far show that both animate and inanimate referents can serve as implied demoted P arguments, although one or the other may be semantically more appropriate depending on the predicate. For example, an animate P argument is implied with predicates such as \textit{-kin} `greet' or \textit{-laɾ} `send/order', while an inanimate participant is more suitable with predicates such as as \textit{-qawalat} `knead' or \textit{-peloɡ} `rake'. Other predicates allow both animate and inanimate P arguments depending on the discourse context, such as \textit{-tʃaq} `cut' or \textit{-lat} `leave'. Finally, note that the distribution of the two agent foregrounding markers cannot be determined by the animacy of the P argument, since both animate and inanimate referents may appear in \textit{-ɢan} and \textit{-taɢan} detransitivized constructions.  
% Mocoví follows this typological trend in that ingestive verbs can be causativized although they are two-place predicates. Conversely, these predicates cannot be detransitivized with the valence modifiers \textit{-ɢan} or \textit{-taɢan}. To illustrate the causativization process via \textit{-ɢan}, consider examples in (\ref{ex:causeeat}) with the verb \textit{-keʔe} `eat$_1$'.  



%\noindent As expected, the verb \textit{-eʔet} `drink' can also be \textit{ɢan-}causativized. In (\ref{ex:causedrinka}), we find the occurrence of that predicate in a transitive clause.  Unfortunately, I did not have the chance to document this verb in its causative version, but is documented in Buckwalter's vocabulary, as illustrated in (\ref{ex:causedrinkb}).


  
      %The only verbal requirement is that the modified stems correspond to a two-place predicate. Second, within the group of two-place predicates, these detransitivizers are much more frequent with those predicates taking the Set II paradigm than Set III. This is why we mostly find \textit{i-} or \textit{∅-} for the third-person subject, instead of \textit{n-}. The latter bound form is the typical third person marker from Set III. %\tinytodo{There is a correlation between antipassives and manner verbs, see Malchukov 2015, Rappaport 2015 and Say 2021. Do I have any regular correlation that support that claim? Mention this?}

%This distributional pattern also corroborates that, regardless of a few exceptions, the predicates derived by \textit{-ɢan} or \textit{-taɢan} will take \textit{ɾ-}, indicating the intransitivity feature associated with this derivation.  Additionally, one notable difference between these two detransitivizers is that only \textit{-ɢan} can also function as a valence increasing marker; \textit{-taɢan} is restricted to the intransitivizing domain.

%\subsection{Complement clauses}
%The foregrounding constructions occur in different syntactic contexts, most of which favor topic. 
% \begin{exe}
% \ex \label{ex: }
% \begin{xlist}
% \ex 
% \gll ʔwe ka i-waːn-aʁan-aʁak ∅-tʃikqoʔot ka ɲ-iwenat-aʁak ɾa i-waːn-aʁan-aʁak  ma̰ʔ qa-ʒ-apaʁaɣin ka s-adenat-\textbf{aʁan} \\
%  \\
% \trans ‘.' \label{ex:}

%\ex 
%\gll \\
%\\
%\trans ‘.'  \label{ex:}
% \end{xlist}
%\end{exe}
 
%\subsection{Lexicon building interaction}{\label{sec:}}
 Another relevant fact about the verb markers \textit{‑ɢan} and \textit{‑taɢan} is their crucial role in the creation of derived lexical items from verbal sources. Morphologically derived nouns of different types often require a change in the valence profile of the base predicate before undergoing category change via nominalizers. The topic of nominalizations in this languages is complex, involving different nominalizerd markers, and will not be treated with much detail here. For illustrative purpose only, I show how the valence-changing operator \textit{-ɢan} works in tandem with what \citet[][150-153]{gualdieri1998} defined as agent nominalization morphology. 
 
    

\begin{exe}
\ex 
\gll ʔwe ka i-waːn-\textbf{ɢan}-ɢak  \\
{\Exist} {\DetOne} {\First\Sg.\Poss}-see-{\Vm:\Intr}-{\Nmlz.\Act}\\
\trans `... (and)  I have a vision (for the future) ...'  \hfill (mocCA190924\_1: 00:05:58) \label{ex:nom-action}
\end{exe}
 
 
 \begin{exe}
\ex 
\gll ka-i-ɾa ɾa i-aʔɢat-\textbf{ɢan}-ɢak na naʔɢa \\
 {\DetOne}-{i}-{\DetFour} {\DetFour} {\First\Sg.\Poss}-speak-{\Vm:\Intr}-{\Nmlz.\Act} {\DetThree} day\\
\trans `...these are my words today...'  \hfill (mocCA190924\_1: 00:07:00) \label{ex:nom-action-two}
\end{exe}
 
 \begin{exe}
\ex 
\gll ʔwe so ɾa l-eːnat-\textbf{ɢan}-ɢak  \\
 {\Exist} {\DetTwo} {\DetThree} {\Third.\Poss}-ask-{\Vm:\Intr}-{\Nmlz.\Act}\\
\trans `...he had questions (today)...'  \hfill (mocCA180709: 00:09:05) \label{ex:nom-action-three}
\end{exe}
 

 Building on the valence restrictions for the selection of \textit{‑ɢan} or \textit{‑taɢan} described above, detransitivized predicates create a new base on which the nominalizer \textit{‑ɢak} operates, as shown in (\ref{ex:nom-action})–(\ref{ex:nom-action-three}). As in typical lexical nominalization processes, the core arguments of the derived predicates are recategorized as participants within the nominal domain. In particular, the subject/agent becomes the possessor of the derived noun. Semantically, it is not arbitrary that \textit{‑ɢan} is selected in the creation of derived nouns of this type, as they denote abstract entities corresponding to the action carried out by the possessor. The nominalization process overrides any implication of the demoted P argument, which explains why the possessor is the only remaining participant. 
 
 %%%% Copy the structure of examples. They should be in this format. 
% \begin{exe}
% \ex \label{ex: }
% \begin{xlist}
% \ex 
% \gll i-paɾaɢan-aɢa-ɾ-pi \\
%  {\First\Sg.\Poss}-hunt-{\Nmlz}-{\Pl}-{\Coll}\\
% \trans ‘my hunting friends.' \label{ex:}

%\ex 
%\gll \\
%\\
%\trans ‘.'  \label{ex:}
% \end{xlist}
%\end{exe}
 

%INCLUDE HERE THE EXAMPLE FROM ASK in Nieves text.


\section{Discourse distribution: a case study}{\label{sec:annotation-scheema}}
One aspect that has not received much attention in relation to agent foregrounding constructions concerns their discourse distribution. This section therefore examines the actual occurrence of \textit{‑ɢan} and \textit{‑taɢan} in natural speech. The goal is purely descriptive: to determine how frequent the different types of agent foregrounding predicates are in discourse, and to identify any correspondence between predicate types and the expression of P arguments  \citep[for a similar example, see][]{haude.allassonniere-tang2023}. Although these questions are basic, such a quantitative investigation has not previously been carried out for valence‑changing operations in Mocoví or in other Guaycuruan languages. There are, however, relevant studies on related topics, such as counts of lexical argument expression in Mocoví \citep[see][]{califa2014} and Chaco Toba's discourse \citep[see][]{cuneo.messineo2019}.

The case study presented here is based on a single narrative discourse, used as a first step to complement the qualitative observations on the distribution and function of agent foregrounding constructions. The narrative has the characteristics of a monologue, in which the speaker recalls various past and present life events. Building on the clause types and patterns of argument expression described in \ref{sec:argument expression}, each clause is annotated according to the following parameters: (i) predicate meaning, (ii) valence type, (iii) P‑argument expression, and (iv) the type of expression used for the P argument. An additional parameter is applied only to third‑person P arguments, identifying the determiner type used in nominal P‑argument expressions. This variable is intended to capture whether some of the semantic properties encoded by nominal determiners show a correspondence with the types of P‑argument expressions attested in discourse.

\begin{figure}[H] 
\centering 
    \includegraphics[width=1\textwidth]{/Users/cristian/Documents/postdocs/research/abst_pprs_docmnt/2025-pdemot-moc/figures/proportion-predicates.png}
    \caption{Proportion of predicates in natural discourse (values on each bar represent individual counts and percentages).}
    \label{fig:predicates-proportion}
\end{figure}

\autoref{fig:predicates-proportion} shows the proportions of predicates according to their valence type. The labels on the x‑axis indicate whether the predicate is morphologically basic or derived. If predicates are double‑derived via \textit{‑ɢan} or \textit{‑taɢan}, the order and function of the derivational elements follow the same order represented in the labels. Thus, a two‑place predicate may be detransitivized by \textit{‑taɢan} and then transitivized again by \textit{‑ɢan}, a sequence captured by the label \textit{two‑place‑taɢan\_intr‑ɢan\_tr}.

As for the the proportion of predicates, basic one-place and two-place predicates are more frequent than the morphologically derived predicates. In contrast, predicates that include \textit{‑ɢan} or \textit{‑taɢan} in their agent foregrounding function occur with low frequency. Similarly, predicates that combine the transitivizing function of \textit{‑ɢan} with spatial morphology -- thus extending the number of participants encoded by the predicate -- also appear infrequently in discourse. \autoref{tab:predicates-derivation} lists all predicates participating in the \textit{‑ɢan} / \textit{‑taɢan} verb alternation, confirming the tokens found in the larger list presented earlier in \autoref{tab:predicates taking -ɢan and -taɢan}.

\begin{table}[H]
\centering
\begin{tabular}{cllc}
  \hline \hline
 count & verb\_in\_sentence & predicate\_type & n \\ 
  \hline 
1 & speak & two-place-ɢan\_intr-ext &   4 \\ 
  2 & know & two-place-taɢan\_intr &   3 \\ 
  3 & work & two-place-ɢan\_intr &   2 \\ 
  4 & ask\_for\_permission & two-place-ɢan\_intr &   1 \\ 
  5 & eat & two-place-ɢan\_tr &   1 \\ 
  6 & eat & two-place-ɢan\_tr-taɢan\_intr &   1 \\ 
  7 & know & two-place-taɢan\_intr-aɢan\_tr &   1 \\ 
  8 & remember & two-place-taɢan\_intr &   1 \\ 
  9 & speak & two-place-ɢan\_intr &   1 \\ 
  10 & transfer\_knowledge & two-place-taɢan\_intr &   1 \\ 
   \hline \hline
\end{tabular}
\caption{Predicates occurring in the \textit{-ɢan} or \textit{-taɢan} verb alternations.} 
\label{tab:predicates-derivation}
\end{table}

As mentioned in previous sections, there is variation in whether the P argument of agent foregrounding constructions is expressed or not. \autoref{tab:p_expression} illustrates this variation. Although the sample is small, the variation is still observable at the discourse level. Note that all instances of derived predicates in which the final derivational element corresponds to the intransitivizing function of \textit{‑ɢan} or \textit{‑taɢan} include cases where the demoted P argument is absent. In contrast, when \textit{‑ɢan} functions as a transitivizing element or when argument‑extension morphology is involved, the P argument is present.

\begin{table}[H]
\centering
\begin{tabular}{clcc}
  \hline \hline
count & predicate\_type & P\_present & P\_absent \\ 
  \hline
1 & two-place-ɢan\_intr-ext &   4 &   0 \\ 
  2 & two-place-taɢan\_intr &   2 &   3 \\ 
  3 & two-place-ɢan\_intr &   1 &   3 \\ 
  4 & two-place-taɢan\_intr-aɢan\_tr &   1 &   0 \\ 
  5 & two-place-ɢan\_tr &   1 &   0 \\ 
  6 & two-place-ɢan\_tr-taɢan\_intr &   0 &   1 \\ 
   \hline \hline
\end{tabular}
\caption{Occurrence of presence and absence of P arguments in \textit{-ɢan} or \textit{-taɢan} derived predicates.} 
\label{tab:p_expression}
\end{table}

Looking at the three instances in which agent foregrounding constructions co‑occur with a P argument, we can observe that the discourse contexts involve predicates that highlight the agent’s activity, a reading reinforced by the low referentiality of the demoted argument. In (\ref{ex:p-expression-indef}), for example, the detransitivized predicate is followed by an object complement‑like clause whose predication can be interpreted as an abstract reference to the types of things a teacher teaches. Similarly, (\ref{ex:p-expression-indef-two}) involves an indefinite nominal P argument within a nominal structure that is preceded by the indefinite nominal \textit{naʔwe} ‘all’.

\begin{exe}
\ex 
	 \gll s-aʔden-q-ta so-ʔe maestɾo{\Sploan}  ɾ-apaɢin-\textbf{taɢan}  \textbf{ka} \textbf{sa}-\textbf{s}-\textbf{aʔden}-\textbf{q}     qomiɾ i-aʔdeːn-taɢan-ɢan-eɾ  \\
			         {\First.II}-know-{\Pl}-{\Dur} {\DetTwo}-{?} teacher {\Third\II.\Intr}-transfer\_knowledge-{\Vm:\Intr} {\DetOne}  {\Neg}{\First.\II}-know-{\Pl} {\First\Pl.\Pron} {\Third.\II}-know-{\Vm:\Intr}-{\Vm:\Tr}-{\Pl}\\
	\trans `...we know (what a) teacher (does), he teaches what we don't know, (he) makes us to know (or he instructs us)...'
  \hfill(mocCA190924\_1: 00:04:27)\label{ex:p-expression-indef}
\end{exe}


\begin{exe}
\ex 
\gll s-oʔwenat-\textbf{aɢan}-tak  \textbf{naʔwe}-\textbf{na}-\textbf{he} \textbf{n}-\textbf{owenat}-\textbf{ek}   \\
{\First.\II}-work-{\Vm:\Intr}-{\Prog} {all}-{\DetFive}-{?} {\Indet.\Poss}-work-{\Nmlz.\Obj.\Sg}\\
\trans `...(now) I keep working on any kind of jobs...' \hfill(mocCA190924\_1: 00:03:48)\label{ex:p-expression-indef-two}
\end{exe}


\section{Conclusions and outlook}{\label{sec:conclusions}}
This study has analyzed agent foregrounding constructions in Northern Chaco Mocoví and their effects on the expression of core arguments in basic and derived two‑place predicates. It argues that agent foregrounding constructions, morphologically marked by \textit{-ɢan} or \textit{-taɢan}, primarily indicate the high prominence of the agent‑like argument, while secondarily affecting the expression of the P argument, which may be omitted, expressed as a bare noun, or marked as an oblique. Semantically, these strategies of P‑argument demotion typically involve low definiteness and low discourse referentiality of the demoted nominal referents. Importantly, the paper also shows that the language employs multiverbal constructions to reintroduce the originally demoted P argument by repeating verb roots with identical semantics. This type of multiverbal strategy does not appear to be widespread cross‑linguistically and therefore warrants further typological research.

The lexical distribution of \textit{‑ɢan} and \textit{‑taɢan} shows a complementary and asymmetrical pattern: most predicates participating in agent foregrounding are two‑place verbs, and \textit{‑ɢan} is considerably more frequent than \textit{‑taɢan}. These observations align with typological work on antipassive‑like constructions and illustrate the semantic constraints that shape valence alternation in the language.

The study also highlights the importance of agent foregrounding in lexicon building. The detransitivized predicate provides the base for several derived nouns, particularly those denoting abstract entities associated with an agent’s activity. The nominalization process incorporates the agent as the possessor while eliminating reference to the demoted P argument, showing how valence alternation interacts with lexical derivation.

Finally, the discourse‑based analysis shows that agent foregrounding constructions are relatively infrequent but occur in specific functional contexts. When they appear, they contribute to the structuring of narrative information, especially by emphasizing the agent’s action and reducing the referential prominence of the P argument. These findings open a promising line of research for future work on valence‑changing mechanisms in Mocoví and related Guaycuruan languages, as well as for broader typological investigations of antipassives and argument demotion.

%This study has shown that agent foregrounding constructions in Northern Chaco Mocoví constitute a crucial domain where valence alternation, argument expression, and discourse structure interact. Through the use of the derivational markers ‑ɢan and ‑taɢan, the language systematically increases the discursive prominence of the agent while demoting the P argument. The demoted P argument may remain unexpressed, be reintroduced through multiverbal constructions, or appear as an oblique‑marked participant—each strategy responding to the need to maintain referential clarity as discourse unfolds.
%A key finding is that the detransitivized predicates require a non‑definite P argument, which constrains both nominal expression and the acceptability of certain referential contexts. The complementary distribution of ‑ɢan and ‑taɢan further indicates a lexical partition of verbal roots involved in agent foregrounding, matching typological observations about antipassive markers cross‑linguistically.
%The broader implications of these constructions extend beyond clause‑level grammar. Agent foregrounding plays a significant role in lexicon building, as detransitivized predicates regularly serve as bases for nominalization processes yielding agentive or action‑related nouns. The study’s small discourse‑based sample confirms that although agent foregrounding constructions are relatively infrequent, their distribution is systematic and functionally meaningful.
%Overall, the findings demonstrate that Mocoví exhibits a typologically rare configuration of valence alternation, where antipassive‑like and causative functions coexist within the same morphological markers and participate in a broader system linking argument structure, discourse organization, and lexical derivation 

\section{Supplementary material}{\label{sec:supplementary}}
\begin{itemize}
\item The complete project including the raw data, annotation and code for data processing is in the following public repository \url{https://github.com/JuarezRC/p-demotion-project.git}

\item  The natural discourse analyzed in this work is part of the open access audiovisual collection \textit{Documentation and Description of Northern Chaco Mocoví} \citep{Juarez2019}. Access to the raw audio, video and pictures of the recording session are available here \url{https://www.elararchive.org/uncategorized/SO_9538d856-a984-491d-87c1-86ab9ed2ce74/}
\end{itemize} 
%%%%%%%%GLOSSES%%%%
\printglossaries

\bibliographystyle{unified.linguistic}
\bibliography{/Users/cristian/Documents/bibliotek/ZotMendlybrary.bib}
%\bibliographystyle{chicago}
\end{document}
